\documentclass[journal]{IEEEtran}
\usepackage{amsmath,amssymb,booktabs,tabularx,array,cite}
\usepackage[hidelinks,hypertexnames=false]{hyperref}
\usepackage{xurl,float}
\hypersetup{pdftitle={Event-Conditioned Certified Decision Reuse on HAPT},pdfauthor={Nguyen Tam Duc}}
\newcolumntype{Y}{>{\raggedright\arraybackslash}X}
\newcommand{\code}[1]{\nolinkurl{#1}}
\newcommand{\R}{\mathbb R}
\title{Event-Conditioned Certified Decision Reuse for Streaming Sensor Inference: A Four-Week Study of SmallCNN1D on HAPT}
\author{Nguyen Tam Duc}
\begin{document}

\maketitle
\begin{abstract}
We study whether a frozen twelve-class human activity classifier can skip inference on overlapping sensor windows while preserving its decisions and reducing CPU latency. Weeks 1--4 established a frozen reference, rebuilt the chronological HAPT dataset, derived a sufficient reuse condition, and evaluated 15 of 17 registered methods on 5,053 previously inspected test windows. The regional methods recorded no useful certified reuse. Exact incremental computation reused intermediate features on 5,034 within-recording pairs, but mean query latency increased from 0.3420 to 1.9273\,ms (+463.47\%). On the paired queries that actually reused features, M03 added 1.5861\,ms per query (+463.86\%). No executed configuration demonstrated both useful certified decision reuse and net acceleration. A new stride-one audit found 97.97\% neighboring top-1 agreement over 322,856 pairs, but no positive eleven-margin IBP estimates among the 49,677 pairs admitted by a diagnostic radius-0.1 box.
\end{abstract}
\begin{IEEEkeywords}
Human activity recognition, runtime verification, interval bound propagation, streaming inference.
\end{IEEEkeywords}
\section{Week 1: Research Question and Reference}
\IEEEPARstart{T}{he} research question is whether a verified region around an evaluated input can permit enough future queries to skip the reference forward pass to offset verification and runtime overhead. Consecutive windows share observations at shifted indices. This motivates reuse without implying a small coordinate distance, equal labels, or an unchanged model decision. The study asks which domain contracts preserve decisions, how often they apply, and whether the resulting system is faster under matched accounting. Its target is agreement with a frozen classifier. Classification accuracy is a separate outcome.
\subsection{Prior Work and Scope}
HAPT provides the inertial recordings and annotated posture transitions used in this study \cite{reyesortiz2015hapt,reyes2016transition}. NoScope illustrates workload-specific inference filtering \cite{kang2017noscope}. DeltaCNN motivates reusing intermediate computation when inputs arrive sequentially \cite{parger2022deltacnn}. Neither plays the role of a matched empirical baseline here. Interval bound propagation (IBP) and CROWN/automatic LiRPA provide background on bounding neural network outputs \cite{gowal2018ibp,zhang2018crown,xu2020automatic}. The project compares these ideas by reuse target, admissible domain, strength of guarantee, and accounting scope. The comparison defines the research question without establishing priority or novelty against an exhaustive corpus.
\subsection{Architecture and Baseline Lock}
The reference is \code{SmallCNN1D} in evaluation mode. Three same-padded temporal convolutions map channels $6\to32\to64\to64$ with kernels $5,3,3$, respectively. Each convolution has batch normalization and ReLU. The first two stages have width-two max pooling. The final $64\times32$ features are flattened and passed through a 128-unit dense layer, ReLU, and a twelve-logit layer. Frozen batch normalization is affine. The reference has 283,692 trainable parameters. The first dense layer has 262,272 parameters including biases (92.45\% of the total). Parameter share measures storage and does not measure saved operations or runtime.
\par
Temporal weight sharing captures local multichannel patterns without assigning a separate kernel to every position. The nominal final convolutional receptive field is 20 samples, with a jump of four samples. Flatten gives the classifier access to feature position, a design choice motivated by six directional transition classes. Global average pooling is invariant to permutations of its input feature positions. Prior convolutions may encode directional patterns, and the project has no matched twelve-class head ablation at this checkpoint. Historical valid-convolution/GAP derivations have different architectural assumptions. The finite feedforward graph offers a concrete verification target. Recurrent and attention verifiers were outside the matched Week 1 comparison.
\par
The Week 1 closeout records a validation-selected seed-43 checkpoint, train-fit normalization, full pipeline execution and an independent baseline artifact audit (65/65 checks). On the 3,572 supervised test windows the selected model recorded 89.05\% accuracy and 0.7763 macro-F1. This establishes an exploratory reference with limited transition support. Under the original Week 1 hypothesis IDs, H1 had descriptive support (2,560/2,617 stable pairs versus 23/44 transition pairs), H3 stated the margin implication, and H2/H4 remained proposed. These IDs predate the H01--H24 ledger. Week 1 left G2 on hold and verification and net cost open \cite{localw1}.
\section{Week 2: Data, Model, and Baseline Audit}
HAPT has 61 recordings from 30 subjects. Each window contains 128 samples at 50\,Hz ($2.56$\,s) from three accelerometer and three gyroscope channels. Stride 64 gives 50\% overlap, with no windows across recording boundaries. For a recording of length $T\ge128$, the window count is $1+\lfloor(T-128)/64\rfloor$. Week 2 regenerated all 17,453 windows and manifest fields and verified 17,392 within-recording aligned pairs across splits. A historical indexing discrepancy was documented. A one-sample annotation-origin change alters 98 assignments, so upstream origin remains a provenance qualification \cite{localweek2}.
\begin{table}[t]
\caption{Canonical HAPT partitions}
\centering
\begin{tabular}{lrrr}\toprule
Split & Subjects & Stream & Supervised\\\midrule
Train & 17 & 9,978 & 6,955\\
Validation & 4 & 2,422 & 1,751\\
Test & 9 & 5,053 & 3,572\\\bottomrule
\end{tabular}
\end{table}
The full reconstruction has 10,361 pure, 1,917 accepted majority, 1,588 ambiguous, and 3,587 unlabeled windows. Pure and accepted majority windows supply supervised targets. Ambiguous and unlabeled windows remain in the stream for agreement and cost analysis. The test stream spans 19 recordings and contains 5,034 valid adjacent pairs. The frozen predictions were reproduced on all 5,053 rows. The audit retained class support and showed that aggregate sample count hides sparse validation transition coverage. Overlap induces dependence, so window-level counts are descriptive and do not provide independent subject-level replicates.
\par
Write $X\in\R^{6\times128}$ for a raw window and $z=N(X)$ for channelwise normalization using train statistics. The full reference is $F(X)=f_\theta(N(X))$ in real arithmetic and $\widehat F$ denotes its frozen float32 execution. Reuse identities include checkpoint, normalization, class order, graph, and numerical semantics. Week 2 preserves the data and model identities used throughout Week 3 and the Week 4 replay.
\section{Week 3: Mathematical and Numerical Foundations}
\subsection{Sufficient Condition for Reuse}
Let $c$ be the class of a trusted anchor already evaluated by the reference. For a declared region $S$ in the stated coordinate system, suppose that the verifier returns
\[
 \ell_{cj}\le \inf_{X\in S}\{F_c(X)-F_j(X)\}
 \qquad (j\ne c).
\]
If $\min_{j\ne c}\ell_{cj}>0$, then every $X\in S$ has $F_c(X)>F_j(X)$ for every competitor, so its unique winner is $c$. With twelve classes this requires eleven valid bounds over the same region. A nonpositive bound is inconclusive. Tie behavior does not affect a strict-margin result. The executable decision requires a justified enclosure of the frozen backend or an explicit valid bridge from the real bounds.
\subsection{Bounds and Domains}
For affine $Wv+b$ over a box with midpoint $m$ and radius $r$, IBP has midpoint $Wm+b$ and radius $|W|r$. ReLU and max pooling are coordinatewise monotone and propagate endpoint bounds. For a margin, one can subtract the corresponding output rows before bounding the final affine map. This may retain cancellation that would be lost by bounding two logits independently. A Lipschitz lower bound is $g_{cj}(a)-L_{cj}r$ on a ball of radius $r$ if $L_{cj}$ is a valid upper bound for the margin slope in that norm. CROWN uses affine relaxations instead of propagating only coordinate intervals.
\par
For an affine lower bound $\alpha^\top z+\beta$ on $z=a+Bu$, $\|u\|_\infty\le1$, minimization gives $\alpha^\top a+\beta-\|B^\top\alpha\|_1$. A diagonal box, a wave box with fixed overlap, and a correlated affine region have distinct geometry. A coordinate hull is conservative when containment is proved and can lose correlations. PCA scores alone do not control orthogonal residuals. Region selection may use observed stream state. Membership is a separate test and uses the same domain as the certificate.
\subsection{Counterexamples and Evidence Boundaries}
Shared sensor samples occur at different positions in consecutive windows. A ramp gives exact overlap and nonzero positional distance. Equal ground-truth labels also fail to imply equal predicted classes, as the Week 4 hypothesis ledger documents. Conversely, continuity and a strict winner at an anchor imply a local neighborhood where the winner is unchanged. This existence result gives no useful lower bound on its size. A compact set strictly inside a continuous winner region has a positive true minimum margin. Loose bounds can still fail to certify it.
\par
For sound executable reuse, the implementation must contain input conversion, normalization, affine reductions, nonlinear operators and final comparisons under a stated arithmetic model. If an executable per-margin error bound is $\eta_{cj}$, requiring $\ell_{cj}>\eta_{cj}$ suffices with conservative membership. Neither an arbitrary tolerance nor a sampled agreement estimate supplies $\eta_{cj}$. Fixed-anchor lifetime and finite-horizon containment must hold throughout the declared horizon. Per-step similarity does not authorize chaining certificates.
\subsection{Alignment, Lifecycle, and Cost}
For the frozen graph the last convolutional feature at index $j$ depends nominally on input indices $[4j-8,4j+11]$. At stride 64, next-window interior columns $j=2,\ldots,13$ align with old columns $j+16$. Boundary and pool alignment conditions matter. Aligned feature reuse still requires the position-sensitive head and does not imply top-1 decision reuse.
\par
A certificate can be reused while its anchor, model, domain and numerical identity remain valid. A hit does not itself justify recentering. Failed membership, invalid identity, inconclusive bounds, recording boundaries, or nonfinite values require fallback or input error handling. A whole-path cost model includes selection, membership, fallback, certificate builds and initialization. With mean overhead $T_o$, reference cost $T_f$, accepted reuse fraction $r$, builds per query $q$, and average build cost $T_b$, the simplified break-even condition is $rT_f>T_o+qT_b$. Timed traces take precedence over this approximation.
\par
Week 3 used validation-selected radius $10^{-6}$ in normalized coordinates. It obtained 5,034 diagonal and 4,211 wave conditional-positive bounds over 5,053 anchors. Neither admitted a next window. It recorded zero formal runtime reuse. Ordinary float64 propagation did not establish executable float32 soundness. The complete path ranged from 3.107 to 3.198\,ms per query against 0.215 to 0.224\,ms for the reference core in that run. Native auto-LiRPA compatibility and surrogate scope were evaluated separately. The Week 3 run and math report retain the counterexamples and method probes \cite{localw3}.
\section{Week 4: Implementation and Replay}
Week 4 moved from conditional estimates to an instrumented comparison: run identity and protocol sealing, numerical primitives and normalization, direct margins, domain and solver controls, incremental inference, runtime and state ownership, fault checks, comparator audits, and full replay. At a recording boundary the runtime resets state. The reference path owns anchor creation. Validation configuration and method execution status are separately recorded. The numerical contract scopes guarantees to supported operations and a declared backend. A fixed tolerance alone does not enclose accumulated floating-point error. An unsupported solver or missing numerical justification cannot authorize a reuse. Recorded method targets and actual execution capabilities are separate.
\par
The preflight specifies the original checkpoint. The saved run uses an exploratory test stream previously inspected in Weeks 1--3. The seal documents the replay protocol without creating a new holdout. The registry contains 17 entries, of which 15 ran and two comparator entries were audited as unsupported. The run record identifies Linux x86-64, PyTorch 2.14.1 CPU, and one Torch thread for the timed benchmark. The reference core (M00), wrapper (M01), and owned reference (M16) have different accounting scopes. M03 reuses intermediate computations. M05--M13 attempted regional methods or associated controls. The supplementary tables reproduce method outcomes and all 24 hypotheses from the run ledger \cite{localw4}.
\par
\subsection{End-to-End Query Latency and Actual Reuse Cost}
The CPU benchmark timed the complete method path from an in-memory query to its decision at batch size one, in three repetitions of 5,053 windows. The timer includes work performed inside the method, such as inference, reused feature processing, and verification or fallback when invoked. Cold setup, warmup, shadow evaluation, trace writing, and harness work outside the query call are outside this scope. M00 mean/p50/p95 were 0.3420/0.3303/0.4070\,ms. We paired each timing observation with the same M00 row and repetition \cite{localpaired}. For a given subset, with baseline mean $B$ and method mean $T$, the latency difference is $\Delta=T-B$ and the percentage reduction is $100(B-T)/B$. Negative reductions correspond to added delay. The comparison matrix lists replay-wide results for every executed method.
\begin{table}[t]
\caption{Paired query latency in ms. Positive $\Delta$ is added time.}
\label{tab:latency-paired}
\centering\small
\begin{tabular}{@{}lrrrr@{}}\toprule
Path & $B$ & $T$ & $\Delta$ & Reduction\\\midrule
M03 all & 0.3420 & 1.9273 & +1.5852 & $-463.47\%$\\
M03 reuse & 0.3419 & 1.9280 & +1.5861 & $-463.86\%$\\
M03 nonreuse & 0.3676 & 1.7192 & +1.3517 & $-367.75\%$\\
M11 no-hit & 0.3420 & 0.4479 & +0.1059 & $-30.96\%$\\
M06 no-hit & 0.3420 & 7.1278 & +6.7858 & $-1983.96\%$\\
\bottomrule\end{tabular}
\end{table}
\par
M03 recorded feature reuse on 5,034 of 5,053 queries (99.62\% of all windows, 100\% of within-recording adjacent pairs). Across the three timing repetitions, this gives 15,102 reuse observations and 57 nonreuse observations. On the matched reuse rows, M03 took 1.9280\,ms versus 0.3419\,ms for direct inference, about $5.64\times$ as long. This is a measured feature reuse path, which still runs the position-sensitive classification head.
\par
M05--M13 had no useful hits. Their no-hit paths averaged 0.4479--7.1278\,ms, adding 0.1059--6.7858\,ms per query (+30.96\% to +1983.96\%). A successful regional hit has no measured latency in this replay. If future measurements give hit latency $H$ and miss latency $F$ under the same scope, break-even requires $rH+(1-r)F<B$. For $H<B<F$, the required hit fraction exceeds $(F-B)/(F-H)$. With M11's observed $F=0.4479$\,ms, the optimistic case $H=0$ would still require more than 23.64\% hits if the miss cost stayed the same. This is a sensitivity calculation. The executed methods agreed with the reference on the full stream, and H24 was not established.
\section{Post-Week 4: Follow-up Evidence}
The transition-support study used span-anchored extraction and focal training on the same subject partition with distinct data policies. Its three-seed macro-F1 was $0.8537\pm0.0020$ on the expanded 4,544-window test policy and $0.8018\pm0.0045$ on the canonical 3,572-window policy. The reported spread is across seeds. The later checkpoint closes the recorded G2-A/G2-B criteria for its policy. It does not replace the frozen replay checkpoint or establish faster ECDR inference \cite{localg2}.
\par
A 1,000-pair feasibility study recorded median normalized $\ell_2$ displacement 25.49 and median Lipschitz radius estimate 0.0149. The median of paired ratios was about 1,563. These estimates help explain the small coverage of tested regions. They do not characterize every possible domain or isolate a biomechanical cause. A separate forward-only CPU scaling benchmark measured 0.1404\,ms for SmallCNN1D and 3.3823\,ms for the Transformer example. A proxy of 1.3346\,ms from the CNN workload yields a modeled break-even near 39.5\% accepted reuse and a modeled ceiling near $2.53\times$ at complete reuse. The interpolated crossover near 27.6 million MACs depends on that proxy and the sampled models. The study also measured a convolution ring-buffer microbenchmark, in which direct convolution remained faster. The convolution subtest reported 0.03421\,ms for the allocating path, 0.03282\,ms for the ring buffer, and 0.01233\,ms for the direct operator. Native Transformer verifier latency and complete stream acceleration remain open \cite{localfeasibility}.
\section{Week 5: Evidence Audit and New Experiments}\label{sec:follow-up}
\subsection{Stride Sensitivity on the Frozen Reference}
We reconstructed the raw six-channel test recordings and formed a complete window grid at strides 1, 4, 16, 32, and 64, starting at offset zero in each recording. The stride-64 windows matched the frozen tensor element for element. The experiment reused the frozen weights and normalization, took top-1 predictions in CPU batches of 256, and compared neighboring windows within each grid. No labels were used for candidate membership. The pair counts differ with stride, so differences are descriptive comparisons of each grid. The input hashes and environment are recorded in the experiment receipt \cite{localstride}.
\begin{table}[t]
\caption{Exploratory stride audit across 19 test recordings. Distances use normalized coordinates. The final column counts candidate membership in a box of radius 0.1.}
\label{tab:stride}
\centering\small
\begin{tabular}{rrrrr}\toprule
Stride & Pairs & Median $\ell_2$ & Agree (\%) & Members\\\midrule
1 & 322,856 & 6.34 & 97.97 & 49,677\\
4 & 80,708 & 15.25 & 96.56 & 5,122\\
16 & 20,171 & 23.85 & 90.87 & 804\\
32 & 10,081 & 25.97 & 85.88 & 327\\
64 & 5,034 & 25.55 & 79.10 & 137\\\bottomrule
\end{tabular}
\end{table}
At every tested stride, no next window was inside a box of radius $10^{-6}$, $10^{-4}$, or $10^{-2}$ around its predecessor. At stride one, the radius-0.1 box admitted 49,677 candidates (15.39\%), seven of which disagreed with the predecessor in the batch predictions. We ran direct all-competitor IBP on every one of these candidates' anchor boxes in ordinary float64, reusing the Week 3 layer rules. Zero returned all eleven positive lower-margin estimates. The IBP pass is a numerical diagnostic, with no validated float32 enclosure or complete runtime timing. Its positive-set agreement is undefined. These results do not substantiate the draft's 87.31\% certified reuse claim. Other regions and policies remain outside this experiment.
\subsection{Replay Recount and Remaining Research Questions}
The new recount matched all 5,053 saved M00 top-1 decisions to the Week 1 baseline by row ID. It confirmed 15 executed methods, 5,034 incremental feature reuse flags for M03, and zero reuse flags for the tested regional methods and controls M05--M13 \cite{localrecount}.
\par
The source review did not establish the draft's $6.31\times$ Transformer claim. A native benchmark needs a validated attention verifier, a trained reference, and a matched stream protocol. The existing forward benchmark supports only a cost sensitivity analysis. The report makes no comparative GAP accuracy claim. Testing such a claim would require matched head training and transition metrics under the same split and label policy. The measured CPU latency does not identify cache residency or isolate biomechanics. Those mechanisms remain open to targeted tests.
\section{Discussion and Conclusion}
The four-week study established a reproducible experimental reference, a mathematical condition for decision reuse, and a negative feasibility result at the canonical stride. The replay distinguishes three obligations: preserving decisions, admitting useful inputs, and saving total runtime. These obligations must be established on the same workload and numerical contract for a positive result. The corpus is one inertial dataset with known test exposure. The measurements concern the recorded host and protocols. They do not establish device energy savings, deployment results or generalization across sensors.
\onecolumn
\raggedbottom
\appendices
\section{Goals and Reproduction}
The following source map identifies the owner of each stage. Historical and explanatory documents provide derivations or design context. Executed result claims come from the corresponding receipts and tables.
\begin{table}[H]
\caption{Evidence ownership by stage}
\centering\small
\begin{tabularx}{\columnwidth}{@{}lY@{}}\toprule
Stage & Primary record\\\midrule
Week 1 & \code{00_OVERVIEW/week1}: brief, requirements, evidence register, hypothesis ledger and closeout.\\
Week 2 & \code{01_DATA}, \code{02_MODEL} and \code{output/week2_foundations_20261103}: reconstruction, support, policy, graph and prediction audit.\\
Week 3 & \code{03_MATH} and \code{output/week3_certification_20261103}: derivations, domain contract, solver probes and execution records.\\
Week 4 & \code{04_SYNTHESIS/week4} and \code{output/week4/week4_initial_20261103}: goals, protocol, numerical reviews, registry, replay, hypotheses and independent reconciliation.\\
Follow-up & \code{output/post_week4}: distinct classification support and scaling studies. \code{04_SYNTHESIS/week5/evidence}: new stride diagnostics.\\
\bottomrule\end{tabularx}
\end{table}
The project goal register distinguishes research obligations from positive outcomes. G0 has a scoped literature baseline. G1 has reconstructed data with provenance caveats. The original G2 is conditional and the expanded G2-A/B study passes its recorded criteria. G3 has a sufficient theorem and fail-closed implementation evidence, with numerical validity scoped to the declared profiles and zero useful regional hits. G4 did not establish useful regional structure in replay. G5 did not establish net speedup. G6 has the editable report, compiled PDF, regeneration scripts, and run receipts. Scientific conditions that remain open are identified in the corresponding sections.
\par
Sources and scripts are rooted at the research checkout. Use the Week 1 source of truth and Week 2 audits for identity, the Week 3 run record for diagnostic evidence, and the Week 4 protocol, traces, comparison matrix and independent audit for replay. The frozen checkpoint SHA-256 is
\code{8ffc5dd8eb0e9cdf832160d074f3ae816ae6cba02adf7214a63a8b623303adb1}.
The train-fit preprocessing identifier is
\code{aa8e1cc20829307187e0a78200f9957f7dd068b1dc26063c2720ad832bb2fd6e}.
The later checkpoint is identified in the G2 receipt. The Week 4 run is \code{week4_initial_20261103}. Its closeout records 40 task receipts, which are separate from regression test counts.
\begin{table}[H]
\caption{Week 4 goals, recorded deliverables and scientific disposition}
\centering\small
\begin{tabularx}{\textwidth}{@{}lYY@{}}\toprule
Goal & Recorded deliverable & Evidentiary scope\\\midrule
G01 & Identity manifest, seal and protected artifact record & Frozen replay. Prior exposure remains disclosed.\\
G02 & Numerical derivation, normalization and backend bridge review & Profile-specific assumptions must hold for execution.\\
G03 & Eleven-margin methods and operator checks & Unsupported adapters cannot authorize reuse.\\
G04 & Diagonal, fixed-overlap and affine domain controls & Candidate membership and useful certified reuse are separate outcomes.\\
G05 & Incremental schedule, equivalence checks and full head & Finite M03 reuse and agreement with negative net savings.\\
G06 & Runtime, state lifecycle, failure paths and cost instrumentation & Fallback evidence on the recorded stream.\\
G07 & Hypothesis ledger, comparator and calibration records & Two comparator entries were not executed.\\
G08 & Full stream, timing and memory artifacts & Matched core comparison produced negative results.\\
G09 & Independent audit and regenerated comparison & Recounts agreed with the saved trace.\\
G10 & Editable and compiled technical reports, with delivery receipt & This synthesis adds the chronology and stride audit.\\
\bottomrule\end{tabularx}
\end{table}
\section{Replay Method Inventory}
\begin{table}[H]
\caption{Week 4 comparison matrix. Times and added delay $\Delta$ are in milliseconds. Reduction is relative to M00. Hits retain each method's semantics. A dash marks an unexecuted comparator.}
\centering\small
\begin{tabularx}{\textwidth}{@{}lYrrrrrrrr@{}}\toprule
ID & Method & Hits & Mean & p50 & p95 & QPS & Speed/M00 & $\Delta$ & Reduction (\%) \\
\midrule
M00 & Frozen full core & 0 & 0.342 & 0.3303 & 0.407 & 2923.71 & 1.0 & +0.00 & +0.00 \\
M01 & Full equivalent wrapper & 0 & 1.6766 & 1.6034 & 2.1147 & 596.44 & 0.204 & +1.33 & -390.19 \\
M02 & Exact input replay & 0 & 0.4788 & 0.451 & 0.6365 & 2088.63 & 0.7144 & +0.14 & -39.98 \\
M03 & Exact incremental computation & 5034 & 1.9273 & 1.7548 & 2.8766 & 518.87 & 0.1775 & +1.59 & -463.47 \\
M04 & Heuristic reuse & 0 & 0.4974 & 0.4664 & 0.6773 & 2010.61 & 0.6877 & +0.16 & -45.41 \\
M05 & Global Lipschitz cache & 0 & 0.5094 & 0.4808 & 0.6781 & 1963.01 & 0.6714 & +0.17 & -48.94 \\
M06 & IBP diagonal box & 0 & 7.1278 & 6.5475 & 10.564 & 140.3 & 0.048 & +6.79 & -1983.96 \\
M07 & IBP fixed-overlap wave & 0 & 6.5641 & 6.1424 & 9.7216 & 152.34 & 0.0521 & +6.22 & -1819.16 \\
M08 & Correlated candidate via coordinate hull & 0 & 0.4626 & 0.4493 & 0.5274 & 2161.59 & 0.7393 & +0.12 & -35.26 \\
M09 & Native correlated bounds & 0 & 0.4481 & 0.4449 & 0.4667 & 2231.51 & 0.7632 & +0.11 & -31.02 \\
M10 & CROWN diagonal & 0 & 0.45 & 0.4489 & 0.4705 & 2222.03 & 0.76 & +0.11 & -31.58 \\
M11 & CROWN wave/correlated & 0 & 0.4479 & 0.4448 & 0.4678 & 2232.56 & 0.7636 & +0.11 & -30.96 \\
M12 & Fixed-anchor policy control & 0 & 0.4966 & 0.47 & 0.5024 & 2013.7 & 0.6887 & +0.15 & -45.19 \\
M13 & Finite-horizon policy control & 0 & 6.1458 & 5.9997 & 8.1202 & 162.71 & 0.0557 & +5.80 & -1696.85 \\
M14 & Faithful LipCache comparator & -- & -- & -- & -- & -- & -- & -- & -- \\
M15 & Prior-work applicability only & -- & -- & -- & -- & -- & -- & -- & -- \\
M16 & OWNED\_REFERENCE baseline & 0 & 0.4674 & 0.4513 & 0.5645 & 2139.37 & 0.7317 & +0.13 & -36.66 \\
\bottomrule\end{tabularx}
\end{table}
\section{Hypothesis Dispositions}
The statements and verdicts below follow the machine-readable Week 4 ledger. The underlying files retain Week 3 dispositions, governing tasks, and finite witnesses.
\begin{table}[H]
\caption{Week 4 ledger, H01--H12.}
\centering\small
\begin{tabularx}{\textwidth}{@{}lYY@{}}\toprule
ID & Statement & Recorded Verdict \\
\midrule
H01 & Shift-aligned overlap is exact across adjacent chronological windows & VERIFIED FINITE \\
H02 & 50\% window overlap implies small ordinary coordinate distance & FALSIFIED UNCONDITIONAL \\
H03 & Same activity annotation implies identical CNN prediction & FALSIFIED UNCONDITIONAL \\
H04 & Strict winner margin and continuity establish a local invariant decision region & DERIVED \\
H05 & Positive valid whole-domain margins are sufficient for decision reuse & DERIVED REAL ARITHMETIC \\
H06 & Sound Lipschitz bound defines certified decision ball & CONDITIONAL THEOREM \\
H07 & IBP certifies every invariant domain & FALSIFIED COMPLETENESS \\
H08 & Direct 11-margin output bounds eliminate intermediate logit dependency loss & DERIVED \\
H09 & CROWN is always cheaper and tighter than alternatives & NOT ESTABLISHED \\
H10 & Temporal/PCA directions align with decision margin normals & NOT ESTABLISHED \\
H11 & Structured domains provide superior certified coverage over matched diagonal boxes & NOT ESTABLISHED \\
H12 & Fixed overlap reduces free coordinates from CW=768 to Cs=384 & DERIVED \\
\bottomrule\end{tabularx}
\end{table}
\begin{table}[H]
\caption{Week 4 ledger, H13--H24.}
\centering\small
\begin{tabularx}{\textwidth}{@{}lYY@{}}\toprule
ID & Statement & Recorded Verdict \\
\midrule
H13 & Actual exact intermediate feature reuse preserves bitwise classifier output & VERIFIED FINITE \\
H14 & Incremental computation yields net latency savings on SmallCNN1D CPU inference & FALSIFIED UNCONDITIONAL \\
H15 & Re-centering after certified reuse hit requires no new numerical proof & FALSIFIED UNCONDITIONAL \\
H16 & Fixed-anchor membership preserves repeated soundness across queries & CONDITIONAL THEOREM \\
H17 & Multi-step finite horizon reachable sets preserve certified invariance & NOT ESTABLISHED \\
H18 & Exact box membership worst-case complexity is Theta(m) in coordinate-access model & DERIVED \\
H19 & Mathematical soundness implies useful streaming coverage & FALSIFIED UNCONDITIONAL \\
H20 & Useful coverage implies net end-to-end speedup & FALSIFIED UNCONDITIONAL \\
H21 & ECDR outperforms CNN baseline or LipCache on HAPT continuous streaming & NOT ESTABLISHED \\
H22 & Real arithmetic proofs automatically transfer soundly to Float32 runtime & FALSIFIED UNCONDITIONAL \\
H23 & Correct temporal event name alone is sufficient for sound decision reuse & FALSIFIED UNCONDITIONAL \\
H24 & There exists an operating regime simultaneously sound, useful, and low-cost on HAPT & NOT ESTABLISHED \\
\bottomrule\end{tabularx}
\end{table}
\section{Week 4 Task Register}
The closeout records all 40 listed entries as verified engineering tasks. These dispositions attest to the scoped deliverables and do not imply positive scientific results.
\begin{table}[H]
\caption{Week 4 engineering tasks by closeout milestone.}
\centering\small
\begin{tabularx}{\columnwidth}{@{}llY@{}}\toprule
Group & Task range & Scope \\
\midrule
M0 & 001--004 & Governance, integrity, specification, environment \\
M1 & 005--010 & Numerical contract, margins, IBP, Lipschitz, CROWN, review \\
M2 & 011--018 & Domains, incremental computation, event triggers, lookup, integration \\
M3 & 019--026 & Runtime, fallback, cache, telemetry, reference, registry, assurance, harness \\
M4 & 027--031 & Workload, baselines, calibration, synthetic probes, protocol seal \\
M5 & 032--034 & Replay, timing, memory and scaling \\
M6 & 035--036 & Comparison and independent audit \\
M7 & 037--040 & Scientific and explanatory reports, rendering, delivery verification \\
\bottomrule\end{tabularx}
\end{table}

\section{Week 1: Reference Selection and Requirement Evidence}
\begin{table}[H]
\caption{Stored validation selection across three training seeds.}
\centering\small
\begin{tabularx}{\textwidth}{@{}Yrrr@{}}\toprule
Fit & Seed & Best epoch & Validation macro-F1\\
\midrule
Candidate & 42 & 14 & 0.836829\\
Selected & 43 & 8 & 0.846785\\
Candidate & 44 & 9 & 0.836958\\
\bottomrule\end{tabularx}
\end{table}
Training used a maximum of 30 epochs, batch size 128, Adam at \(10^{-3}\), weight decay \(10^{-4}\), cosine scheduling and patience 10. The selection criterion was validation macro-F1. The baseline test result uses the selected seed and the frozen supervised denominator of 3,572. The evidence locators below follow the Week 1 requirements and closeout \cite{localw1}.
\begin{table}[H]
\caption{Week 1 requirement-to-evidence map.}
\centering\small
\begin{tabularx}{\textwidth}{@{}lYY@{}}\toprule
ID & Requirement & Recorded evidence\\
\midrule
W1-01 & Research grounding & Brief comparison and primary bibliography\\
W1-02 & K=12 lock & selection.json and SOURCE\_OF\_TRUTH.md give the selected reference and input identity\\
W1-03 & Architecture & Executed parameter count and layer shapes, reproduced in Week 2\\
W1-04 & Pipeline execution & Week 1 train, evaluation and verification receipts\\
W1-05 & Transition support & G2 hold and support evidence in the brief\\
W1-06 & Research question & Decision preservation, useful coverage and total cost\\
W1-07 & Evidence register & EVIDENCE\_REGISTER.md and independent baseline check (65/65)\\
W1-08 & Sensor/data contract & Recording, channel order, window grid, split and label policy\\
W1-09 & Hypotheses & H1 to H4 in the Week 1 ledger\\
W1-10 & Stratification & Descriptive stable/transition pair counts and reference agreement\\
W1-11 & Roadmap & Closeout describes handoff to reconstruction and verification\\
W1-12 & English deliverable & Editable Week 1 brief and study records\\
\bottomrule\end{tabularx}
\end{table}
\textbf{Interpretation.} Seed 43 led the three candidates on validation. High agreement for stable adjacent pairs coexists with much lower agreement for transition pairs. The baseline identifies the decision target and motivates the distinction between temporal similarity and a certificate.
\section{Week 2: Dataset and Model Results}
Raw inputs are ordered accelerometer \(x,y,z\), then gyroscope \(x,y,z\). The accepted majority threshold is 77 of 128 samples. Pure and accepted majority windows enter supervised metrics, while ambiguous and unlabeled windows retain their chronological positions. The local reconstruction establishes the project grid and labels. The upstream annotation-origin sensitivity remains a limitation for comparisons that relabel boundary windows \cite{localweek2}.
\begin{table}[H]
\caption{Frozen class order, supervised support and selected-reference test metrics.}
\centering\small
\begin{tabularx}{\textwidth}{@{}rYrrrrrr@{}}\toprule
ID & Class & Train & Val. & Test & Precision & Recall & F1\\
\midrule
1 & Walking & 1077 & 246 & 531 & 0.9603 & 0.9567 & 0.9585\\
2 & Walking Upstairs & 973 & 247 & 537 & 0.9331 & 0.9088 & 0.9208\\
3 & Walking Downstairs & 894 & 231 & 492 & 0.9237 & 0.9837 & 0.9528\\
4 & Sitting & 1103 & 289 & 546 & 0.7997 & 0.8407 & 0.8196\\
5 & Standing & 1206 & 312 & 597 & 0.8798 & 0.8090 & 0.8429\\
6 & Laying & 1188 & 316 & 586 & 0.9831 & 0.9949 & 0.9890\\
7 & Stand To Sit & 79 & 19 & 45 & 0.7619 & 0.7111 & 0.7356\\
8 & Sit To Stand & 53 & 13 & 32 & 0.7778 & 0.6562 & 0.7119\\
9 & Sit To Lie & 90 & 26 & 51 & 0.5893 & 0.6471 & 0.6168\\
10 & Lie To Sit & 89 & 15 & 46 & 0.7714 & 0.5870 & 0.6667\\
11 & Stand To Lie & 119 & 18 & 68 & 0.5714 & 0.5294 & 0.5496\\
12 & Lie To Stand & 84 & 19 & 41 & 0.4737 & 0.6585 & 0.5510\\
\bottomrule\end{tabularx}
\end{table}
\begin{table}[H]
\caption{Channel normalization fitted only on 6,955 supervised training windows.}
\centering\small
\begin{tabularx}{\textwidth}{@{}Yrr@{}}\toprule
Channel & Mean & Std. dev.\\
\midrule
Acc X & 0.8043609 & 0.4057797\\
Acc Y & 0.0411366 & 0.4071547\\
Acc Z & 0.0817700 & 0.3549119\\
Gyro X & 0.0091356 & 0.4476351\\
Gyro Y & -0.0043976 & 0.4212706\\
Gyro Z & -0.0071832 & 0.2964415\\
\bottomrule\end{tabularx}
\end{table}
\textbf{Interpretation.} Transition validation counts are much smaller than basic-activity counts. The lower transition F1 scores identify where the reference is less accurate on this split. They measure agreement with activity labels. Reuse is evaluated against the reference's decisions, including its mistakes.
\begin{table}[H]
\caption{Actual module shapes at batch size one and audited cost per sample.}
\centering\small
\begin{tabularx}{\textwidth}{@{}lYYYrr@{}}\toprule
Module & Operator & Input shape & Output shape & Parameters & MACs\\
\midrule
conv.0 & Conv1d & [1, 6, 128] & [1, 32, 128] & 992 & 122880\\
conv.1 & BatchNorm1d & [1, 32, 128] & [1, 32, 128] & 64 & 0\\
conv.2 & ReLU & [1, 32, 128] & [1, 32, 128] & 0 & 0\\
conv.3 & MaxPool1d & [1, 32, 128] & [1, 32, 64] & 0 & 0\\
conv.4 & Conv1d & [1, 32, 64] & [1, 64, 64] & 6208 & 393216\\
conv.5 & BatchNorm1d & [1, 64, 64] & [1, 64, 64] & 128 & 0\\
conv.6 & ReLU & [1, 64, 64] & [1, 64, 64] & 0 & 0\\
conv.7 & MaxPool1d & [1, 64, 64] & [1, 64, 32] & 0 & 0\\
conv.8 & Conv1d & [1, 64, 32] & [1, 64, 32] & 12352 & 393216\\
conv.9 & BatchNorm1d & [1, 64, 32] & [1, 64, 32] & 128 & 0\\
conv.10 & ReLU & [1, 64, 32] & [1, 64, 32] & 0 & 0\\
head.0 & Flatten & [1, 64, 32] & [1, 2048] & 0 & 0\\
head.1 & Linear & [1, 2048] & [1, 128] & 262272 & 262144\\
head.2 & ReLU & [1, 128] & [1, 128] & 0 & 0\\
head.3 & Linear & [1, 128] & [1, 12] & 1548 & 1536\\
\bottomrule\end{tabularx}
\end{table}
Total cost is 1,172,992 convolution/dense MACs per sample with the archived convention. Bias, normalization, nonlinear operators and memory traffic are outside that count. The Week 2 mean/p50/p95 forward-only measurements were 0.1892/0.1557/0.3111\,ms on a resident normalized input. These measurements do not include the complete Week 4 decision path.
\begin{table}[H]
\caption{Confusion counts on 3,572 supervised windows. Rows are true labels, columns are predictions; class IDs follow the support table.}
\centering\small
\begin{tabularx}{\textwidth}{@{}Yrrrrrrrrrrrr@{}}\toprule
True/Pred. & 1 & 2 & 3 & 4 & 5 & 6 & 7 & 8 & 9 & 10 & 11 & 12\\
\midrule
1 & 508 & 0 & 22 & 0 & 0 & 0 & 0 & 0 & 0 & 0 & 0 & 1\\
2 & 16 & 488 & 15 & 0 & 0 & 0 & 2 & 0 & 0 & 3 & 2 & 11\\
3 & 4 & 2 & 484 & 0 & 0 & 0 & 2 & 0 & 0 & 0 & 0 & 0\\
4 & 0 & 21 & 0 & 459 & 57 & 0 & 1 & 2 & 1 & 0 & 5 & 0\\
5 & 0 & 3 & 0 & 106 & 483 & 0 & 1 & 0 & 1 & 0 & 2 & 1\\
6 & 0 & 0 & 0 & 0 & 0 & 583 & 0 & 0 & 0 & 0 & 3 & 0\\
7 & 1 & 2 & 0 & 4 & 1 & 0 & 32 & 3 & 1 & 0 & 1 & 0\\
8 & 0 & 2 & 0 & 1 & 7 & 0 & 1 & 21 & 0 & 0 & 0 & 0\\
9 & 0 & 0 & 0 & 0 & 1 & 6 & 0 & 0 & 33 & 0 & 11 & 0\\
10 & 0 & 1 & 0 & 2 & 0 & 0 & 0 & 0 & 0 & 27 & 1 & 15\\
11 & 0 & 1 & 1 & 2 & 0 & 4 & 2 & 0 & 20 & 0 & 36 & 2\\
12 & 0 & 3 & 2 & 0 & 0 & 0 & 1 & 1 & 0 & 5 & 2 & 27\\
\bottomrule\end{tabularx}
\end{table}
The full-stream top-two gap had minimum 0.0001135, median 5.9782 and p95 11.8876 logit units across 5,053 inputs. A sampled near-tie motivates strict margins. A large point gap alone gives no lower bound over a region.
\section{Week 3: Domain, Backend and Runtime Evidence}
\begin{table}[H]
\caption{Offline domain diagnostic at the validation-selected normalized radius 1e-6.}
\centering\small
\begin{tabularx}{\textwidth}{@{}Yrrrrr@{}}\toprule
Domain & Builds & Positive estimates & Next members & Formal hits & Mean build (ms)\\
\midrule
diagonal box & 5053 & 5034 & 0 & 0 & 6.774\\
fixed overlap wave & 5053 & 4211 & 0 & 0 & 6.830\\
\bottomrule\end{tabularx}
\end{table}
\begin{table}[H]
\caption{Week 3 operational replay. Each run used 5,053 windows and zero accepted reuse.}
\centering\small
\begin{tabularx}{\textwidth}{@{}Yrrrrrr@{}}\toprule
Domain & Repeat & Builds & Fallback & Core (ms) & ECDR (ms) & Reduction (\%)\\
\midrule
diagonal box & 1 & 49 & 5053 & 0.224 & 3.174 & -1316.92\\
fixed overlap wave & 1 & 69 & 5053 & 0.224 & 3.198 & -1327.35\\
diagonal box & 2 & 49 & 5053 & 0.215 & 3.107 & -1346.81\\
fixed overlap wave & 2 & 69 & 5053 & 0.215 & 3.169 & -1375.60\\
\bottomrule\end{tabularx}
\end{table}
\begin{table}[H]
\caption{Fallback reason counts in the timed Week 3 runs.}
\centering\small
\begin{tabularx}{\textwidth}{@{}YrY@{}}\toprule
Domain & Repeat & Reason and query count\\
\midrule
diagonal box & 1 & no certificate: 2, outside domain: 5,004, stale or expired anchor: 29, recording reset: 18\\
fixed overlap wave & 1 & no certificate: 22, outside domain: 4,984, stale or expired anchor: 29, recording reset: 18\\
diagonal box & 2 & no certificate: 2, outside domain: 5,004, stale or expired anchor: 29, recording reset: 18\\
fixed overlap wave & 2 & no certificate: 22, outside domain: 4,984, stale or expired anchor: 29, recording reset: 18\\
\bottomrule\end{tabularx}
\end{table}
\begin{table}[H]
\caption{Saved auto-LiRPA probes on the K=12 graph.}
\centering\small
\begin{tabularx}{\textwidth}{@{}YYrrY@{}}\toprule
Graph & Conversion & Bound records & Positive & Scope\\
\midrule
native 1d & FAILED & 0 & 0 & No native bounds\\
singleton height 2d & PASS & 12 & 6 & Conditional surrogate diagnostics\\
\bottomrule\end{tabularx}
\end{table}
\textbf{Interpretation.} The native 1-D operator conversion did not produce bounds in this probe. The singleton-height surrogate matched sampled logits, but the probe did not establish exact executable equivalence in general. Positive interval estimates remained conditional, and no next query entered the tested domains. Builds, state checks and fallback therefore added cost without avoiding reference forwards \cite{localw3}.
\section{Week 4: Resource Use and Task Evidence}
\begin{table}[H]
\caption{Recorded attributable adapter/cache memory and membership coordinate counts.}
\centering\small
\begin{tabularx}{\textwidth}{@{}Yrrrrr@{}}\toprule
Method & Adapter B & State B & Certificate B & Total B & Coordinates\\
\midrule
M00 & 1113 & 0 & 0 & 1113 & 768\\
M01 & 1307 & 0 & 0 & 1307 & 768\\
M02 & 1533 & 3080 & 0 & 4613 & 768\\
M03 & 12465 & 3584 & 0 & 16049 & 384\\
M04 & 1586 & 6664 & 0 & 8250 & 384\\
M05 & 8909 & 0 & 6144 & 15053 & 768\\
M06 & 11940 & 0 & 6144 & 18084 & 768\\
M07 & 11705 & 0 & 6144 & 17849 & 768\\
M08 & 1460123 & 0 & 6144 & 1466267 & 768\\
M09 & 1460411 & 0 & 6144 & 1466555 & 768\\
M10 & 1459976 & 0 & 6144 & 1466120 & 768\\
M11 & 1460238 & 0 & 6144 & 1466382 & 768\\
M12 & 11769 & 0 & 6144 & 17913 & 768\\
M13 & 11815 & 0 & 6144 & 17959 & 768\\
M14 & -- & -- & -- & -- & --\\
M15 & -- & -- & -- & -- & --\\
M16 & 1514 & 0 & 0 & 1514 & 768\\
\bottomrule\end{tabularx}
\end{table}
The separately reported frozen float32 tensor payload is 1,136,072 bytes. The float64 bounds snapshot is 2,272,120 bytes. The attribution table tracks the fields recorded for each method. Process RSS includes the runtime, backend and resident arrays, so it cannot be interpreted as cache cost. The study did not isolate dynamic workspace or energy \cite{localw4}.
\textbf{Interpretation.} M03 provides actual feature reuse, but its timed hit path is already slower than a complete M00 decision. Regional methods incur costs with no observed useful hits. Low attributed state size does not itself imply low query latency.
\begin{table}[H]
\caption{Week 4 closeout tasks 001--020. Receipt locators are relative to the Week 4 run directory.}
\centering\small
\begin{tabularx}{\textwidth}{@{}lYYl@{}}\toprule
Task & Scope in closeout & Receipt locator & Status\\
\midrule
W4-001 & Task tracker, schemas and lock governance & tasks/W4-001/receipt.json & VERIFIED\\
W4-002 & Baseline snapshot, dirty tree guard and integrity audit & tasks/W4-002/receipt.json & VERIFIED\\
W4-003 & Method taxonomy, protocol matrix and task crosswalk & tasks/W4-003/receipt.json & VERIFIED\\
W4-004 & Runtime, profiling and hardware telemetry & tasks/W4-004/receipt.json & VERIFIED\\
W4-005 & Numerical soundness and floating-point safety bounds & tasks/W4-005/receipt.json & VERIFIED\\
W4-006 & K=12 competing-class margin analysis & tasks/W4-006/receipt.json & VERIFIED\\
W4-007 & IBP engine and radius calculation & tasks/W4-007/receipt.json & VERIFIED\\
W4-008 & Global and local Lipschitz bounds & tasks/W4-008/receipt.json & VERIFIED\\
W4-009 & LP/CROWN adapter & tasks/W4-009/receipt.json & VERIFIED\\
W4-010 & Proof obligations and verification review & tasks/W4-010/receipt.json & VERIFIED\\
W4-011 & Ball regions & tasks/W4-011/receipt.json & VERIFIED\\
W4-012 & Box and diagonal regions & tasks/W4-012/receipt.json & VERIFIED\\
W4-013 & Polytope and zonotope adapters & tasks/W4-013/receipt.json & VERIFIED\\
W4-014 & PCA and low-dimensional projection & tasks/W4-014/receipt.json & VERIFIED\\
W4-015 & Incremental streaming and partial recomputation & tasks/W4-015/receipt.json & VERIFIED\\
W4-016 & Event triggers and heuristic detectors & tasks/W4-016/receipt.json & VERIFIED\\
W4-017 & Exact and near-neighbor temporal lookup & tasks/W4-017/receipt.json & VERIFIED\\
W4-018 & Domain integration and compliance & tasks/W4-018/receipt.json & VERIFIED\\
W4-019 & Runtime adapter and lifecycle & tasks/W4-019/receipt.json & VERIFIED\\
W4-020 & Fail-closed arbitration and fallback & tasks/W4-020/receipt.json & VERIFIED\\
\bottomrule\end{tabularx}
\end{table}
\begin{table}[H]
\caption{Week 4 closeout tasks 021--040. Receipt locators are relative to the Week 4 run directory.}
\centering\small
\begin{tabularx}{\textwidth}{@{}lYYl@{}}\toprule
Task & Scope in closeout & Receipt locator & Status\\
\midrule
W4-021 & Cache and eviction & tasks/W4-021/receipt.json & VERIFIED\\
W4-022 & Query instrumentation & tasks/W4-022/receipt.json & VERIFIED\\
W4-023 & Reference and shadow validation & tasks/W4-023/receipt.json & VERIFIED\\
W4-024 & Registry and lifecycle checks & tasks/W4-024/receipt.json & VERIFIED\\
W4-025 & Fault injection and assurance & tasks/W4-025/receipt.json & VERIFIED\\
W4-026 & Benchmark harness & tasks/W4-026/receipt.json & VERIFIED\\
W4-027 & Workload and split isolation & tasks/W4-027/receipt.json & VERIFIED\\
W4-028 & Reference core and other baseline paths & tasks/W4-028/receipt.json & VERIFIED\\
W4-029 & Calibration and threshold selection & tasks/W4-029/receipt.json & VERIFIED\\
W4-030 & Synthetic perturbation benchmarks & tasks/W4-030/receipt.json & VERIFIED\\
W4-031 & Protocol seal and freeze & tasks/W4-031/receipt.json & VERIFIED\\
W4-032 & Complete stream replay & tasks/W4-032/receipt.json & VERIFIED\\
W4-033 & End-to-end CPU timing & tasks/W4-033/receipt.json & VERIFIED\\
W4-034 & Memory and scaling profiles & tasks/W4-034/receipt.json & VERIFIED\\
W4-035 & Synthesis and comparison & tasks/W4-035/receipt.json & VERIFIED\\
W4-036 & Independent recount and audit & tasks/W4-036/receipt.json & VERIFIED\\
W4-037 & Scientific report & tasks/W4-037/receipt.json & VERIFIED\\
W4-038 & Explanatory report & tasks/W4-038/receipt.json & VERIFIED\\
W4-039 & Rendered report and layout review & tasks/W4-039/receipt.json & VERIFIED\\
W4-040 & Delivery verification and handoff & tasks/W4-040/receipt.json & VERIFIED\\
\bottomrule\end{tabularx}
\end{table}
\section{Week 4: Hypothesis Witnesses}
\begin{table}[H]
\caption{Evidence recorded in the Week 4 hypothesis ledger, H01--H08.}
\centering\small
\begin{tabularx}{\textwidth}{@{}lYYY@{}}\toprule
ID & Recorded verdict & Witness or interpretation & Governing tasks\\
\midrule
H01 & VERIFIED FINITE & 5,034 adjacent within-recording test pairs verified; 12 intermediate Conv1D feature positions bitwise equal. & W4-001, W4-002, W4-007, W4-015, W4-019, W4-027\\
H02 & FALSIFIED UNCONDITIONAL & Ramp shifts and high-frequency sensor perturbations produce large L2/L\_inf coordinate deltas despite 50\% temporal overlap. & W4-002, W4-015, W4-018, W4-027, W4-035\\
H03 & FALSIFIED UNCONDITIONAL & Rows 682 and 683 share true ID 11 (stand-to-lie); zero-based predictions are 10 and 8. & W4-023, W4-027, W4-032, W4-035\\
H04 & DERIVED & Topological continuity theorem guarantees open neighbourhood around query with Delta > 0, but certified radius remains tiny under sensor scale. & W4-005, W4-011, W4-012, W4-027\\
H05 & DERIVED REAL ARITHMETIC & Soundness theorem holds under exact real arithmetic; float32 implementation requires closed numerical bridge (W4-010). & W4-005, W4-006, W4-007, W4-008, W4-009, W4-010, W4-022, W4-036\\
H06 & CONDITIONAL THEOREM & M05 global spectral norm Lipschitz constant yields r = Delta / (2 * L\_global). Non-trivial radius on continuous stream yields R=0 hits. & W4-005, W4-006, W4-007, W4-008, W4-009, W4-010, W4-012, W4-028, W4-033\\
H07 & FALSIFIED COMPLETENESS & IBP interval propagation incurs severe dependency loss across 3 Conv1D layers and MLP head, failing to certify true invariant boxes. & W4-006, W4-008, W4-010, W4-011, W4-027\\
H08 & DERIVED & On the same hidden box, direct margin minimization is no worse than separate extrema. & W4-011, W4-014, W4-018, W4-032, W4-035\\
\bottomrule\end{tabularx}
\end{table}
\begin{table}[H]
\caption{Evidence recorded in the Week 4 hypothesis ledger, H09--H16.}
\centering\small
\begin{tabularx}{\textwidth}{@{}lYYY@{}}\toprule
ID & Recorded verdict & Witness or interpretation & Governing tasks\\
\midrule
H09 & NOT ESTABLISHED & The archived probes do not establish universal superiority in cost or tightness. & W4-013, W4-014, W4-028, W4-033, W4-035\\
H10 & NOT ESTABLISHED & Local geometry probes do not establish alignment across certified regions. & W4-002, W4-016, W4-017, W4-018, W4-030, W4-032, W4-035\\
H11 & NOT ESTABLISHED & Both structured PCA affine hulls and matched diagonal boxes achieve exactly 0 useful hits on full test stream under valid budgets. & W4-015, W4-016, W4-017, W4-018, W4-030, W4-032, W4-035\\
H12 & DERIVED & Stride 64 with window 128 fixes 64 samples (384 coordinates) exactly when adjacent window belongs to same continuous recording. & W4-007, W4-015, W4-018, W4-025, W4-034\\
H13 & VERIFIED FINITE & Bitwise aligned features and finite decision agreement; no universal bitwise logit guarantee. & W4-019, W4-020, W4-021, W4-026, W4-032, W4-033, W4-034\\
H14 & FALSIFIED UNCONDITIONAL & Paired trace: M03 1.9273 ms, M00 0.3420 ms; reduction -463.47 percent. & W4-019, W4-020, W4-021, W4-027, W4-033, W4-035\\
H15 & FALSIFIED UNCONDITIONAL & Re-centering around accepted query without fresh bound computation leads to unbounded drift and catastrophic loss of certificate soundness. & W4-022, W4-023, W4-025, W4-026, W4-027\\
H16 & CONDITIONAL THEOREM & Immutability of anchor query and certificate domain guarantees that containment check contains true output class across lifetime. & W4-010, W4-022, W4-023, W4-025, W4-026, W4-032\\
\bottomrule\end{tabularx}
\end{table}
\begin{table}[H]
\caption{Evidence recorded in the Week 4 hypothesis ledger, H17--H24.}
\centering\small
\begin{tabularx}{\textwidth}{@{}lYYY@{}}\toprule
ID & Recorded verdict & Witness or interpretation & Governing tasks\\
\midrule
H17 & NOT ESTABLISHED & Zero useful hits for the horizon control; containment theorem remains conditional. & W4-015, W4-017, W4-022, W4-025, W4-032, W4-033, W4-034\\
H18 & DERIVED & Proof establishes that verifying containment in non-degenerate box requires reading all m coordinates in worst case; verified by microbenchmarks. & W4-007, W4-016, W4-018, W4-026, W4-033, W4-034\\
H19 & FALSIFIED UNCONDITIONAL & Methods M05-M13 possess mathematically valid bounding algorithms but achieve R = 0 (0\% reuse) on the 5,053-query continuous HAPT test stream. & W4-027, W4-030, W4-032, W4-035, W4-036\\
H20 & FALSIFIED UNCONDITIONAL & M03 achieves 100\% within-recording reuse (5,034 hits) but suffers speed ratio 0.1775x vs M00 due to tensor manipulation overhead exceeding forward execution. & W4-024, W4-028, W4-030, W4-032, W4-033, W4-034, W4-035\\
H21 & NOT ESTABLISHED & M00 is fastest among timed paths. Faithful LipCache was not timed. & W4-028, W4-029, W4-030, W4-031, W4-032, W4-033, W4-034, W4-035, W4-036\\
H22 & FALSIFIED UNCONDITIONAL & A justified error enclosure for the executable backend is required. & W4-005, W4-006, W4-007, W4-008, W4-009, W4-010, W4-021, W4-026, W4-036\\
H23 & FALSIFIED UNCONDITIONAL & Windows sharing identical event names (e.g. WALKING) span multiple distinct feature trajectories and target classes. & W4-015, W4-022, W4-023, W4-026, W4-027\\
H24 & NOT ESTABLISHED & No witness among 15 executed out of 17 registered; other regimes remain open. & W4-010, W4-021, W4-030, W4-031, W4-032, W4-033, W4-034, W4-035, W4-036, W4-037, W4-038, W4-039, W4-040\\
\bottomrule\end{tabularx}
\end{table}
The witnesses have different logical scope. A finite observation applies to the recorded data. A mathematical counterexample rejects a universal assertion. A conditional theorem needs its domain and numerical assumptions before executable reuse can follow.
\section{Follow-up: Transition Quality and Stride Sensitivity}
\begin{table}[H]
\caption{Separate post-Week 4 G2 training. Canonical and expanded test policies use different denominators.}
\centering\small
\begin{tabularx}{\textwidth}{@{}Yrrrrrrr@{}}\toprule
Selection & Seed & Canonical acc. & Canonical F1 & Transition F1 & Expanded acc. & Expanded F1 & Transition F1\\
\midrule
Epoch 9 & 42 & 0.8989 & 0.8065 & 0.6929 & 0.8779 & 0.8519 & 0.7893\\
Epoch 6 & 43 & 0.8956 & 0.7956 & 0.6707 & 0.8779 & 0.8528 & 0.7908\\
Epoch 14 & 44 & 0.9009 & 0.8032 & 0.6832 & 0.8818 & 0.8565 & 0.7951\\
\bottomrule\end{tabularx}
\end{table}
The expanded policy removed all 18 recorded support blockers (18 to zero). Mean canonical test accuracy was 89.85 percent and mean expanded test accuracy was 87.92 percent. Mean transition F1 values were 0.6822 and 0.7917 under those respective policies. The training comparison changes the data policy and objective, so it does not isolate the effect of either change. The frozen Week 4 replay and its timing still refer to the original reference \cite{localg2}.
\begin{table}[H]
\caption{Stride audit. Counts use chronological pairs within each grid and recording.}
\centering\small
\begin{tabularx}{\textwidth}{@{}Yrrrrrr@{}}\toprule
Stride & Pairs & Equal top-1 & Median L2 & Median Linf & Box hits & Hit disagreements\\
\midrule
1 & 322856 & 316298 & 6.342 & 1.420 & 49677 & 7\\
4 & 80708 & 77931 & 15.251 & 2.776 & 5122 & 1\\
16 & 20171 & 18329 & 23.852 & 3.493 & 804 & 0\\
32 & 10081 & 8658 & 25.970 & 3.870 & 327 & 0\\
64 & 5034 & 3982 & 25.546 & 3.934 & 137 & 0\\
\bottomrule\end{tabularx}
\end{table}
\textbf{Interpretation.} Shorter strides increase empirical neighboring agreement under these grids. Seven stride-one pairs within the radius-0.1 box changed decision. All 49,677 radius-0.1 member pairs were checked with conditional IBP and none had eleven positive lower bounds. This explains why high temporal agreement does not substantiate certified reuse for the tested box \cite{localstride}.
\textbf{Research conclusion.} The four-week work established the frozen data/model contract, derived a sufficient decision-preservation rule, and tested multiple ways to reuse work. The measured canonical result is negative for acceleration. M03 demonstrated feature reuse at a latency penalty, and the tested regional paths supplied no useful certified hits. The separate G2 study improved classification under its stated support policy. Future claims of net certified acceleration require a valid executable certificate, actual hits and a measured complete-path advantage on the same workload.

\clearpage
\section{Proofs, Counterexamples, and Evaluation of H01--H24}
The identifiers below follow the Week 4 ledger. Week 3 H14 instead concerns the historical GAP head. Let $W$ be window length, $s$ stride, $C$ channel count, and $g_{cj}=F_c-F_j$ the margin against competitor $j$. The mathematical arguments use real arithmetic unless specified. The empirical conclusions use the recorded configurations and workload \cite{localw3,localw4,localpaired}.

\subsection{Overlap and Certification}
\paragraph{H01. Exact shift-aligned overlap.}
Write $X_i(:,t)=u(:,is+t)$ for one recording $u$. For $0\le t<W-s$,
\[
X_{i+1}(:,t)=u(:,(i+1)s+t)=X_i(:,s+t).
\]
This proves the shifted equality for 64 samples per channel at $W=128$, $s=64$. The trace covers 5,034 same-recording adjacent pairs. No equality across recording boundaries is implied.

\paragraph{H02. Overlap alone does not bound positional distance.}
In one channel take $u(t)=at$. The aligned overlap is exact, but equal-index subtraction gives $as$ at every position. Therefore the $\ell_\infty$ and $\ell_2$ distances are $|a|s$ and $|a|s\sqrt W$. Scaling $a$ makes either arbitrarily large in this construction. Overlap supplies no small-distance guarantee without an additional constraint.

\paragraph{H03. Same annotation can accompany different predictions.}
Frozen manifest rows 682 and 683 are consecutive and both have pure activity ID 11, stand-to-lie. Their zero-based prediction indices are 10 and 8. One example of equal labels and unequal predictions disproves the universal implication. It gives no estimate of the population failure rate.

\paragraph{H04. A strict winner has a local invariant neighborhood.}
Assume every $g_{cj}(a)>0$ and continuity at anchor $a$. For each competitor choose $\delta_j>0$ so that within that distance the margin changes by less than $g_{cj}(a)/2$. The minimum of eleven positive radii is positive. In this common neighborhood all margins stay positive. Existence alone supplies no radius useful on the stream.

\paragraph{H05. Valid positive whole-region margins suffice.}
For every admitted $x\in S$, a valid lower bound gives $g_{cj}(x)\ge\ell_{cj}$. If all eleven $\ell_{cj}>0$, the class wins uniquely everywhere in $S$. For executable reuse, conservative membership must hold and numerical error must be covered. For example, if executable margins differ from the ideal margins by at most $\eta_{cj}$ on $S$, testing $\ell_{cj}>\eta_{cj}$ makes the executable margins positive.

\paragraph{H06. Valid Lipschitz bounds give a certified ball conditionally.}
If each margin is $L_{cj}$-Lipschitz on the proposed ball in the stated norm, then $g_{cj}(x)\ge g_{cj}(a)-L_{cj}\|x-a\|$. For a strict anchor winner, choose $r<\min_{j:L_{cj}>0}g_{cj}(a)/L_{cj}$. Positive constant margins impose no restriction. The proof requires a justified upper bound, not an unchecked norm estimate. M05 recorded zero useful hits.

\subsection{Bounds and Domain Structure}
\paragraph{H07. IBP is incomplete.}
On $t\in[-1,1]$, set $h_1=h_2=\operatorname{ReLU}(t+1)$. The true margin $1+h_1-h_2$ is always 1. IBP replaces the hidden values by the box $[0,2]^2$, giving a lower bound of $1+0-2=-1$. This invariant domain is missed because the box loses the equality constraint. A failed bound does not by itself demonstrate a changed decision.

\paragraph{H08. Direct margins can preserve final-layer cancellation.}
For logits $A,D$ on the same hidden box $B$, every $A(h)-D(h)$ is at least $\inf_B A-\sup_B D$. Taking the infimum proves the direct affine minimum cannot be worse than subtraction of separate extrema. For $A(h)=1+h$ and $D(h)=h$ on $[0,2]$, the direct lower bound is 1 and the separate result is $-1$. This can avoid final-layer dependency loss and cannot restore dependencies already lost inside the graph.

\paragraph{H09. Universal CROWN superiority is unestablished.}
Validity compares a bound with the true minimum and imposes no runtime ordering. For an affine margin on a box, interval minimization is already exact, so a valid relaxation cannot always be strictly tighter. Runtime and practical tightness depend on graph, region, backend and budget. The archived native and surrogate probes do not establish universal dominance.

\paragraph{H10. PCA variance alone does not imply margin alignment.}
For covariance $\operatorname{diag}(4,1)$ and margin $g(x)=1+x_2$, the leading PCA direction is $(1,0)$ and the margin normal is $(0,1)$. Their dot product is zero. The example rules out an implication from variance alone. Local measurements in the frozen CNN do not establish alignment throughout regions where the slopes can change.

\paragraph{H11. Structured coverage did not improve in replay.}
On a common query set, count admitted queries with valid certificates and divide by the query count. The observed numerators for the tested diagonal and structured paths are both zero, yielding zero measured difference. A smaller region might tighten the bound but admit fewer queries. No superiority at another budget follows from this run.

\paragraph{H12. Fixed overlap reduces the number of free coordinates.}
H01 fixes the first $W-s$ samples of each channel of the next window. The tail has $Cs=6\cdot64=384$ free coordinates, compared with $CW=6\cdot128=768$. A tail-only membership check is justified only if the prefix invariant is guaranteed separately.

\subsection{Incremental Inference and Lifecycle}
\paragraph{H13. Aligned features have a derivation and finite validation.}
At position $j$ the final feature has nominal raw support $[4j-8,4j+11]$. Adding stride 64 shifts this to the old position $j+16$. For $j=2,\ldots,13$ both supports are interior and pooling phases coincide. The local operators see the same input samples and fixed normalization, yielding equal real features. The archive records bitwise checks on $5{,}034\cdot12\cdot64=3{,}866{,}112$ corresponding scalars. M03 still reconstructs and evaluates the position-sensitive head, and the recorded replay agrees in top-1. These checks do not prove universal bitwise equality of executable logits.

\paragraph{H14. Incremental computation did not save latency.}
The historical GAP sum can cancel shared contributions. Flatten assigns position-specific weights, so shifting one reused feature can change its contribution. Week 4 H14 tests runtime on the actual graph. Mean latency rises from 0.3420 to 1.9273\,ms, an increase of 1.5852\,ms or 463.47\%. Actual reuse rows add 1.5861\,ms. These measurements reject net savings for this implementation and workload.

\paragraph{H15. A reuse hit does not authorize recentering.}
For $F_1(x)=1-x$, $F_2(x)=0$, class 1 wins on the radius-0.6 interval around zero. The query 0.6 is a valid hit. Reusing the same radius around that query would admit 1.2, where class 2 wins. A certificate for the original interval cannot be transferred without an additional argument.

\paragraph{H16. Fixed-anchor reuse preserves repeated decisions conditionally.}
Keep the domain, certificate and executable identity fixed. Every admitted query passes conservative membership in the same $S$. H05 and the numerical premise apply independently to each query, so every admitted decision is preserved. This gives no guarantee that later queries will belong to $S$.

\paragraph{H17. Horizon preservation is conditional on reachable-set coverage.}
If $R_t$ contains every admissible input at step $t$ and the same class has valid positive margins over $\bigcup_{t=1}^H R_t$ under the numerical contract, H05 applies at every step. The theorem does not guarantee tight reachable sets or useful coverage. The finite-horizon control had zero useful replay hits.

\paragraph{H18. Exact box membership costs $\Theta(m)$ coordinate accesses.}
Checking both bounds for every coordinate costs $O(m)$. If an exact algorithm accepts an interior point after leaving a coordinate unread, change only that coordinate to a value outside its interval. The algorithm sees the same values and accepts, while the correct outcome is rejection. Thus every coordinate must be read on that interior input. The bound concerns coordinate accesses, not measured latency.

\subsection{Coverage, Cost, and Numerical Validity}
\paragraph{H19. Soundness does not imply useful coverage.}
A sound policy states that admission implies correctness. A policy that admits no queries satisfies this implication and has zero coverage, disproving the proposed implication. On the measured regional paths, agreement through fallback likewise supplied zero observed useful hits.

\paragraph{H20. Useful coverage does not imply net speedup.}
With constant hit and miss costs $H,F$, expected cost is $rH+(1-r)F$. Even $r=1$ fails to beat a baseline $B$ if $H>B$. On matched M03 reuse rows the measured hit cost was 1.9280\,ms and M00 cost 0.3419\,ms. This shows high feature reuse with slower execution in the actual workload.

\paragraph{H21. Superiority over M00 or LipCache is unestablished.}
M00 had the lowest complete-path mean, 0.3420\,ms. Every other executed path had a higher mean. M14 did not execute a faithful LipCache comparison, so its matched latency is unavailable. The recorded comparison supplies no acceleration over M00 and does not resolve the external comparison.

\paragraph{H22. A real proof does not automatically prove binary32 behavior.}
Evaluate $g_1=(2^{24}+1)-2^{24}$ in that order and let $g_2=1/2$. Over the reals $g_1=1>g_2$. With binary32 round-to-nearest ties-to-even, $2^{24}+1$ is halfway between representable $2^{24}$ and $2^{24}+2$ and rounds to the former. The subtraction gives zero, reversing the order. A sound transfer requires an error enclosure for the actual numerical operations.

\paragraph{H23. An event name alone is insufficient.}
If two inputs share an event name but their frozen predictions differ, a name-only cache fails when one anchors the other. H03 gives a same-activity example. A name can select a candidate domain. Reuse still requires valid margins and conservative membership under matching identity.

\paragraph{H24. No tested configuration satisfies the joint regime.}
Check decision preservation, useful reuse, and positive net speedup on the same configuration. M03 supplied feature hits and finite decision agreement but failed the speed condition. The regional paths supplied no useful hits or speedup. The baseline paths do not supply useful reuse, and the unsupported entries have no execution evidence. None of the 15 executed configurations out of 17 registered meets the conjunction. Existence over other models or policies remains open.

\bibliographystyle{IEEEtran}
\bibliography{refs}
\end{document}
